\begin{definition}[Norme equivalente]
Soit \todo{Lecture 3} $N_{1}$ et $N_{2}$ deux normes sur $V$. On dit qu'elles sont equivalents s'ils existent $0<c_{1}<c_{2}<\infty$ tels que :
$$
c_{1}N_{1}(\boldsymbol{x})\le N_{2}(\boldsymbol{x})\leq c_{2}N_{1}(\boldsymbol{x}),\quad\forall \boldsymbol{x}\in V
$$\end{definition}
\begin{remark}
On a que les normes "max", euclidienne et $p$ sont equivalents car 
$$
\forall p,q\in [1,\infty)~\exists  \underline{c}_{p,q}, \bar{c}_{p,q}: \underline{c}_{p,q}\lVert \boldsymbol{x} \rVert _{p}\le\lVert \boldsymbol{x} \rVert _{q}\le \bar{c}_{p,q}\lVert \boldsymbol{x} \rVert _{p},\quad\forall \boldsymbol{x}\in \mathbf{R}^{n}
$$
\end{remark}
\begin{proposition}
\emph{Toutes} les normes sur $\mathbf{R}^{n}$ sont équivalentes.
\end{proposition}

\begin{definition}[Distance]
Soit $X$ un ensemble, on definit une distance symetrique sur $X$ comme une application $d:X\times X \to\mathbf{R}$ tel que :
\begin{enumerate}
\item $\forall x,y\in X$ on a $d(x,y)\ge0$ et $d(x,y)=0$ si $x=y$.
\item $\forall x,y\in X$ on a $d(x,y)=d(y,x)$.
\item $\forall x,y,z\in X$ on a $d(x,y)\le d(x,z)+d(z,y)$.
\end{enumerate}
\end{definition}

Etant donne un espace vectoriel reel norme $(V,\lVert \bullet \rVert_{N})$, la norme induit une distance qu'on definit par 
$$
d(\boldsymbol{x},\boldsymbol{y})=\lVert \boldsymbol{y}-\boldsymbol{x} \rVert _{N}
$$

\begin{definition}[Espace metrique]
Un ensemble $X$ avec une distance $d$ est appelle \emph{espace metrique} qu'on note $(X,d)$.
\end{definition}
\begin{remark}
Un espace vectoriel avec produit scalaire implique un espace vectoriel norme qui implique un espace metrique.
\end{remark}
\begin{definition}[Produit scalaire euclidien]
Pour tous $\boldsymbol{x}$,$\boldsymbol{y}$ dans $\mathbf{R}^{n}$ on peut definir:
$$
\langle \boldsymbol{x},\boldsymbol{y} \rangle =\boldsymbol{x}^{\top}\boldsymbol{y}=\boldsymbol{y}^{\top}\boldsymbol{x}=\sum_{i=1}^{n} x_{i}y_{i}
$$
\end{definition}
\begin{definition}
Un produit scalaire sur un espace vectoriel reel $V$ est une application $b:V\times V\to \mathbf{R}$ qui:
\begin{enumerate}
\item $\forall \boldsymbol{x},\boldsymbol{y}\in V$ on a $b(\boldsymbol{x},\boldsymbol{y})=b(\boldsymbol{y},\boldsymbol{x})$.
\item $\forall \boldsymbol{x},\boldsymbol{y},\boldsymbol{z}\in V$ et $\alpha,\beta \in \mathbf{R}$ on a $b(\alpha \boldsymbol{x}+\beta \boldsymbol{y},\boldsymbol{z})=\alpha b(\boldsymbol{x},\boldsymbol{z})+\beta b(\boldsymbol{y},\boldsymbol{z})$.
\item \label{prop3} $\forall \boldsymbol{x}\in V$ on a $b(\boldsymbol{x},\boldsymbol{x})\ge 0$ et $b(\boldsymbol{x},\boldsymbol{x})=0$ si $\boldsymbol{x}=\boldsymbol{0}$.
\end{enumerate}
\end{definition}

\begin{lemma}[Inégalité de Cauchy-Schwarz]
Soit $(V,b)$ un espace vectoriel avec produit scalaire $b$, alors $\forall \boldsymbol{x},\boldsymbol{y}\in V$ on a :
$$
|b(\boldsymbol{x},\boldsymbol{y})|\le\sqrt{ b(\boldsymbol{x},\boldsymbol{x}) }\sqrt{ b(\boldsymbol{y},\boldsymbol{y}) }=\lVert \boldsymbol{x} \rVert _{E}\lVert \boldsymbol{y} \rVert_{E} 
$$
\end{lemma}
\begin{proof}
Pour tous $\boldsymbol{x},\boldsymbol{y}\in V$ et $\alpha \in \mathbf{R}$ on a que 
\begin{align*}
0&\le b(\alpha \boldsymbol{x}+\boldsymbol{y},\alpha \boldsymbol{x}+\boldsymbol{y})=\alpha b(\boldsymbol{x},\alpha \boldsymbol{x}+\boldsymbol{y})+b(\boldsymbol{y},\alpha \boldsymbol{x}+\boldsymbol{y})\\
&=\alpha^{2}b(\boldsymbol{x},\boldsymbol{x})+2\alpha b(\boldsymbol{x},\boldsymbol{y})+b(\boldsymbol{y},\boldsymbol{y}) \\
\Rightarrow\Delta & =4b(\boldsymbol{x},\boldsymbol{y})^{2}-4b(\boldsymbol{x},\boldsymbol{x})b(\boldsymbol{y},\boldsymbol{y})\le 0
\end{align*}

\end{proof}
\begin{lemma}
Soit $(V,b)$ un espace vectoriel reel avec produit scalaire. Alors l'application $\lVert \bullet \rVert_{b}:V\to \mathbf{R}$ donne par 
$$
\lVert \boldsymbol{x} \rVert _{b}=\sqrt{ b(\boldsymbol{x},\boldsymbol{x}) }
$$
est une norme sur $V$.
\end{lemma}
\begin{proof}
On veut demontrer la propriete ~\ref{propiete3}

\begin{align*}
\lVert \boldsymbol{x}+\boldsymbol{y} \rVert ^{2}_{b} & =b(\boldsymbol{x}+\boldsymbol{y},\boldsymbol{x}+\boldsymbol{y}) \\
 & =b(\boldsymbol{x},\boldsymbol{x})+2b(\boldsymbol{x},\boldsymbol{y})+b(\boldsymbol{\boldsymbol{y},\boldsymbol{y}}) \\
 & \leq b(\boldsymbol{x},\boldsymbol{x})+2\sqrt{ b(\boldsymbol{x},\boldsymbol{x}) }\sqrt{ b(\boldsymbol{y},\boldsymbol{y}) }+b(\boldsymbol{y},\boldsymbol{y}) \\
 &\le \left(\sqrt{ b(\boldsymbol{x},\boldsymbol{x}) }+\sqrt{ b(\boldsymbol{y},\boldsymbol{y}) }\right)^{2}=(\lVert \boldsymbol{x} \rVert _{b}+\lVert \boldsymbol{y} \rVert _{b})^{2}
\end{align*}

\end{proof}

\section{Suites dans \texorpdfstring{$\mathbf{R}^{n}$}{Rn}}
\begin{definition}
Soit $\{ \boldsymbol{x}^{(k)} \}_{k\in \mathbf{N}}$ une suite dans $\mathbf{R}^{n}$ qu'on note
$$
\boldsymbol{x}^{(k)}=(x_{1}^{(k)},\dots,x_{n}^{(k)})\in \mathbf{R}^{n}
$$
on dit que $\{ \boldsymbol{x}^{(k)} \}_{k\in \mathbf{N}}$ converge a $\boldsymbol{x}\in \mathbf{R}^{n}$ si
$$
\lim_{ k \to \infty } \lVert \boldsymbol{x}^{(k)}-\boldsymbol{x} \rVert=0
$$
qu'on peut definir par epsilon-delta comme:
$$
\forall\varepsilon>0\exists N_{\varepsilon}>0\forall k\ge N_{\varepsilon}:\lVert \boldsymbol{x}^{(k)}-\boldsymbol{x} \rVert <{\varepsilon}
$$
\end{definition}
\begin{remark}
Cette notion de convergence ne depend pas de la norme choisie. En particulier on peut choisir $\lVert \bullet \rVert_{\infty}$ car si $\lim_{ k \to \infty }\lVert \boldsymbol{x}^{(k)}-\boldsymbol{x} \rVert=0$ alors
$$
\forall\varepsilon>0\exists N_{\varepsilon}>0\forall k\geq N_{\varepsilon}:\max_{i=1,\dots,n}|x^{(k)}_{i}-x_{i}|<\varepsilon
$$
alors si pour tout $i$
$$ 
|x^{(k)}_{i}-x_{i}|<\varepsilon \Rightarrow \lim_{ k \to \infty }x^{(k)}_{i}=x_{i}
$$

Reciproquement: supposons que 
$$
\forall\varepsilon>0\exists N_{\varepsilon}>0\forall k\ge N_{\varepsilon}:|x_{i}^{(k)}-x_{i}|<\varepsilon
$$
Soit $N_{\varepsilon}=\max_{i}N_{\varepsilon}^{(i)}$ alors $\forall k\ge N_{\varepsilon}$ on a que $\max_{i}|x^{(k)}_{i}-x_{i}|<\varepsilon$ donc $\lim_{ k \to \infty }\lVert \boldsymbol{x}^{(k)}-\boldsymbol{x} \rVert=0$
\end{remark}
\begin{definition}[Suite de Cauchy]
On dit qu'une suite $\{ \boldsymbol{x}^{(k)} \}_{k\in \mathbf{N}}\subset \mathbf{R}^{n}$ est de Cauchy si 
$$
\forall\varepsilon>0\exists N_{\varepsilon}>0\forall k,\ell\geq N_{\varepsilon}:\lVert \boldsymbol{x}^{(k)}-\boldsymbol{x}^{(\ell)} \rVert <\varepsilon
$$
\end{definition}
\begin{theorem}
Une suite $\{ \boldsymbol{x}^{(k)} \}_{k\in \mathbf{N}}\subset \mathbf{R}^{n}$ est convergente si et seulement si elle est de Cauchy.
\end{theorem}
\begin{proof}
Une suite $\{ \boldsymbol{x}^{(k)} \}$ converge a $\boldsymbol{x}$ si et seulement si $x_{i}^{(k)}\to x_{i}$ pour tout $i$ si et seulement si $\{ \boldsymbol{x}_{i}^{(k)} \}$ est de Cauchy si et seulement si $\{ \boldsymbol{x}^{(k)} \}$ est de Cauchy.
\end{proof}
\begin{definition}
Une suite $\{ \boldsymbol{x}^{(k)} \}_{k\in \mathbf{N}}\subset \mathbf{R}^{n}$ est bornee s'il existe $M>0$ tel que
$$
\lVert \boldsymbol{x}^{(k)} \rVert\leq M\quad\forall k
$$
\end{definition}
\begin{theorem}[Bolzano-Weierstrass sur $\mathbf{R}^n$]\label{BW}
Etant donnee une suite $\{ \boldsymbol{x}^{(k)} \}_{k\in \mathbf{N}}\subset \mathbf{R}^{n}$ bornee, on peut toujours extraire une sous-suite $\{ \boldsymbol{x}^{(k_{i})} \}_{i\in \mathbf{N}}$ convergente.
\end{theorem}
\begin{proof}
Soit une suite $\{ \boldsymbol{x}^{(k)} \}_{k\in \mathbf{N}}\subset \mathbf{R}^{n}$. Pour tout $\boldsymbol{x}^{(k)}$ on peut regarder la $i$-eme composante et extraire une sous-suite convergente pour ce sous-indice $x^{(k)}_{i}$. On procede en parcourant toutes les elements de la suite en prenant compte a chaque tour de seulement compter le sous-suites que sont convergentes avec la suite du element precedent pour a la fin avoir un indice de sous-suites $(k_{i})_{i\in \mathbf{N}}$ tel que toutes les $x^{(k)}_{i}$ convergent.
\end{proof}